\chapter{Giới thiệu đề tài}

\section{Bối cảnh thực tiễn}

Thực tiễn hiện nay cho thấy nhu cầu hỗ trợ về chăm sóc sức khỏe ban đầu và tiếp cận thông tin y tế của người dân ngày càng tăng cao. Tuy nhiên, quá trình này đang đối mặt với một số thách thức cụ thể như sau:

\begin{itemize}
    \item Người dùng có nhu cầu được hỗ trợ nhanh chóng trong việc nhận diện, đánh giá các triệu chứng phổ biến như cảm cúm, dị ứng, rối loạn tiêu hóa... Tuy nhiên, việc tiếp cận các dịch vụ tư vấn y tế kịp thời và chính xác còn hạn chế, đặc biệt ở khu vực ngoài đô thị hoặc trong các tình huống khẩn cấp.
    
    \item Nhận thức và kiến thức về chăm sóc sức khỏe tại nhà của phần lớn người dân còn hạn chế. Nhiều người chưa nắm vững các biện pháp xử lý ban đầu đối với các vấn đề sức khỏe thông thường, dẫn đến nguy cơ tự ý điều trị sai cách hoặc bỏ qua các dấu hiệu cảnh báo quan trọng.
    
    \item Đa số người dùng không chuyên gặp khó khăn trong việc tiếp cận, tra cứu và hiểu các nguồn thông tin y khoa chính thống. Điều này làm gia tăng nguy cơ tiếp nhận thông tin sai lệch hoặc áp dụng các biện pháp chăm sóc sức khỏe không phù hợp, ảnh hưởng tiêu cực đến hiệu quả phòng bệnh và điều trị.
\end{itemize}

\section{Bối cảnh công nghệ}

Bên cạnh các thách thức thực tiễn, sự phát triển nhanh chóng của trí tuệ nhân tạo (AI) trong lĩnh vực y tế, đặc biệt là các hệ thống đa tác nhân (multi-agent) và AI đa phương thức (multimodal AI), đang mở ra nhiều cơ hội mới. Các mô hình này cho phép hệ thống hiểu, kết hợp và xử lý đồng thời nhiều dạng dữ liệu như văn bản, hình ảnh, âm thanh... Nhờ đó, khả năng nhận diện, phân tích triệu chứng và cung cấp thông tin y khoa ngày càng chính xác, cá nhân hóa và phù hợp với từng người dùng. Tuy nhiên, việc tích hợp hiệu quả các công nghệ này vào thực tiễn vẫn là một bài toán lớn, đòi hỏi giải pháp tổng thể và linh hoạt.

Tổng quan, bối cảnh thực tế đặt ra yêu cầu cấp thiết đối với việc phát triển các giải pháp hỗ trợ tiếp cận thông tin y tế chính xác, dễ hiểu cho mọi đối tượng người dùng. Việc giải quyết các thách thức này có ý nghĩa quan trọng trong việc nâng cao hiệu quả chăm sóc sức khỏe ban đầu, giảm tải cho hệ thống y tế và góp phần bảo vệ sức khỏe cộng đồng một cách bền vững.

\section{Mục tiêu}

Từ những phân tích trên, đề tài hướng tới việc xây dựng một hệ thống Multi-Agent trợ lý y tế ảo có khả năng phân tích triệu chứng đa phương thức và cung cấp thông tin y khoa đáng tin cậy cho người dùng phổ thông.

Việc ứng dụng mô hình đa tác nhân kết hợp AI đa phương thức sẽ giúp hệ thống tiếp nhận, xử lý và giải thích nhiều loại dữ liệu khác nhau, từ đó hỗ trợ người dùng hiệu quả hơn trong việc nhận diện và xử lý các vấn đề sức khỏe thường gặp. 